\section{Introduction}\label{sec:intro}

Stock futures on the National Stock Exchange of India settle at the volume-weighted average
price (VWAP) of the underlying share in the cash market between \DesignWindowStart{} and
\DesignWindowEnd{} on the last Thursday of the month. The settlement price is therefore
produced by the cash market, and anyone carrying a large expiring position has a direct
financial interest in what that market does for \DesignWindowMinutes{} minutes.

This is a well-known design trade-off. A settlement struck at a single closing print is
cheap to move, so exchanges generally do not use one. An average over a window is more
expensive to move, because the mover has to keep trading for the length of the window, and
each additional unit of size has less effect on the average than the last. The question is
empirical: does the window look different, and if so, in what way.

The answer is not obvious in either direction, because expiry brings a second kind of
activity that has nothing to do with the settlement price. Positions have to be closed or
rolled, hedges have to be unwound, and that alone would make the window busier, wider and
harder to trade. Any comparison against a normal day therefore measures both effects
together. Every number in this report is a difference against a matched control session in
the same month, which removes the market environment but not this second effect.

\subsection*{What order-level data adds}

Most studies of settlement windows work from trade prints and top-of-book quotes. Those show
what happened but not what was attempted. An order entered and cancelled a fraction of a
second later leaves no trace in a trade file. A large order shown a few hundred shares at a
time looks, in a depth feed, like a stream of unrelated small orders.

The exchange publishes order-by-order files that record each entry, modification and
cancellation as a separate event, with the participant category and the total and displayed
quantities attached. Three of the measures in this report cannot be built from trade data at
all: the rate at which orders are cancelled within a second of entry, the share of resting
size that is never displayed, and the rate at which concealed size is replenished as it
fills. Rebuilding the book from these events also gives the displayed size at each price
level, which is what makes it possible to ask how much an order of a given size would move
the price.

\subsection*{The question this report tries to answer}

Almost everything one can measure in a busy window---wider spreads, more cancellations,
larger price impact---rises when there is simply more trading, regardless of why. Reporting
those alone cannot separate a market being pushed from a market that is merely crowded.

We therefore build the report around measures that respond to the two explanations with
opposite signs. If prices are being pushed in one direction, returns lean the same way from
one second to the next and the variance ratio rises above one; if the window is busier and
more two-sided, the extra activity arrives as bid-ask bounce and the ratio falls. If one
participant is working a large order, signed order flow is persistent; if many participants
are trading against each other, it is not. More trading on its own cannot move either
statistic in the direction that indicates pressure.

We also compute the settlement price. The original design for this project measured many
properties of the window and never calculated the number the window exists to produce, which
left the central question---does the settlement end up somewhere unusual---unasked.
Section~\ref{sec:results} reports how far it sits from the mid at the window's open, from
the closing price, and from the final minute's own average.

Section~\ref{sec:background} describes the settlement mechanism.
Section~\ref{sec:data} covers the data and the book reconstruction, Section~\ref{sec:method}
the measures and the tests, and Section~\ref{sec:results} the findings.
Section~\ref{sec:discussion} reads them together, and Section~\ref{sec:limitations} sets out
what this design cannot establish.
